\documentclass[11pt]{article}
\usepackage[T1]{fontenc}
\usepackage[margin=1in]{geometry}
\usepackage{amsmath,amssymb,amsthm,array,longtable}
\usepackage[hidelinks]{hyperref}
\newcommand{\Q}{\mathbb Q}
\newcommand{\C}{\mathbb C}
\newcommand{\Z}{\mathbb Z}
\newcommand{\kk}{{\overline{\mathbb Q}}}
\DeclareMathOperator{\CH}{CH}
\DeclareMathOperator{\Nm}{Nm}
\DeclareMathOperator{\Gal}{Gal}
\DeclareMathOperator{\Hdg}{Hdg}
\DeclareMathOperator{\Sp}{Sp}
\DeclareMathOperator{\SO}{SO}
\DeclareMathOperator{\SL}{SL}
\DeclareMathOperator{\im}{im}
\DeclareMathOperator{\Tr}{Tr}
\newtheorem{theorem}{Theorem}
\newtheorem{lemma}[theorem]{Lemma}
\newtheorem{proposition}[theorem]{Proposition}
\newtheorem{corollary}[theorem]{Corollary}
\theoremstyle{remark}
\newtheorem{remark}[theorem]{Remark}
\title{All-degree CM supply for quadratic actions:\\
a constructive verification of the determinant reduction}
\author{Campaign 9000032: Engineer submission}
\date{November 5, 2026}
\begin{document}
\maketitle

\begin{abstract}
We verify the algebraic reduction in the two all-powers theorems of
\emph{Weil classes and Hodge classes on abelian powers}
(October 6, 2026). Reusing the independently reviewed actual derived-group
theorem, we prove that every rational Hodge class on every self-power in
either family is supplied by the full cohomology of CM abelian varieties
through rational algebraic correspondences. The proof covers arbitrary
tensor degrees and repeated determinant factors. It includes a fixed-list
simultaneous-balance argument, primitive integer exponent cancellation,
nonzero determinant extraction, actual positive polarization padding, and
rational Hodge lifts. Markman's all-period Weil-plane theorem, not the
manuscript's later geometric proof, is an external input.
An executable exact certificate checks the finite case reduction and
symbolic identities; it does not replace the universal proof.
The final Hodge-conjecture implication remains conditional on the separate
CM Hodge theorem, which is not audited here.
\end{abstract}

\noindent\textbf{Status and comparison.}
This is a universal prose verification with an exact algebraic certificate,
submitted for independent Reviewer certification. It is not a Lean/kernel
proof, a counterexample, or a novelty certification.
The prior reviewed result \texttt{AG-UNIVERSAL} established condition (c)
of the manuscript's tensor criterion, namely the actual independent
derived images. The present result establishes the additional
determinant-supply condition (d) and the all-degree passage to rational
classes. It does not assume independence of central characters, nor does
it replace an actual special-period Hodge group by a generic PEL group.

\medskip
\noindent\textbf{Source scope.}
The relevant staged labels are \texttt{p3:cm-operations},
\texttt{p3:admissible-operations}, \texttt{p3:tensor-criterion},
\texttt{p3ab:hyperbolic}, and \texttt{p3ab:determinant-supply},
together with the balance arguments at source lines 1609--1834.
For admissible operations we prove the abelian-variety instance used
in these two families. Algebraicity assertions in the CM-operations
lemma and tensor criterion retain their explicitly separate CM input.

\section{Statement, scope, and inputs}

Fix $\kk=\overline{\Q}\subset\C$. Scalar conjugation acts on coefficients
of the cohomology of a \emph{fixed} complex variety. It does not conjugate
that variety. For a smooth projective $X$ define
\[
 {\cal C}^{d}(X)=
 \sum_{\substack{B\ {\rm CM}\\ j\in\Z,\ \Gamma}}
 \im\left(\Gamma_*:H^{d-2j}(B,\Q)(-j)
                 \longrightarrow H^d(X,\Q)\right),\qquad
 \Gamma\in\CH^{\dim B+j}(B\times X)_\Q.
 \tag{1}\label{eq:supply}
\]
Only existing degrees and cycle groups are used; a point is allowed.
CM means that the degree-one Hodge group is a torus.
The twist convention is $V(r)^{p,q}=V^{p+r,q+r}$.
A $\kk$-subspace is supplied if it lies in ${\cal C}^d(X)_\kk$.
In particular, the source in (1) is \emph{full cohomology}, not just
its Hodge classes.

\begin{theorem}\label{thm:main}
Let $A/\C$ admit a unital action of an imaginary quadratic field $E$.
Suppose either
\begin{enumerate}
\item $\dim A\leq5$; or
\item $\dim A=6$ and a compatible polarization has hyperbolic rational
homological $E$-Hermitian form of signature $(3,3)$.
\end{enumerate}
For every $m\geq1$ and $p\geq0$,
\[
 H^{2p}(A^m,\Q)\cap H^{p,p}(A^m)
             \ \subset\ {\cal C}^{2p}(A^m).
 \tag{2}\label{eq:main}
\]
Every such class has a rational Hodge lift through a finite sum of
the defining maps in (1).
Consequently, \emph{if} the rational Hodge conjecture holds for every
CM abelian variety, it holds for every $A^m$ in the two families.
\end{theorem}

We reuse the reviewed theorem in
\path{research/actual_group_audit/report.tex}, without reclassifying groups.
After isogeny and admissible matrix-copy identifications, it supplies the
following covering list. Here $B$ is a simple non-CM factor, $g=\dim B$,
$r$ its multiplicity, and a Type IV center has degree $2f$.
\begin{center}\small
\begin{tabular}{p{.17\textwidth}p{.73\textwidth}}
Types I, II & Symplectic columns of ranks $2,4,6$.\\
Type III & $g=4$ or $6$, $r=1$, one $\SO_g$ column;
each $E$-eigenspace is an admissible copy, of signature
$(g/2,g/2)$.\\
Type IV, $E\not\hookrightarrow K$ &
$K$ quadratic, $g=3,r=2$, one $\SL_3$ column and its opposite.\\
Type IV, $E\hookrightarrow K$ &
$K$ is the full endomorphism field. Either $f=1$, $3\leq g\leq6$,
or $(f,a,g)=(2,2,4),(2,3,6),(3,2,6)$, where $a=g/f$.
Each column has image $\SL_a$; opposite columns are dual.
Rank-two cases have both pure and mixed embedding pairs.
\end{tabular}
\end{center}
The corresponding classical factors act independently on their own
columns and trivially on the CM summand. All copy maps and pair tensors
are admissible. A pure column in a non-CM factor is \emph{not} moved into
the rational CM summand. The list is necessary, not a realizability claim.

The following external propositions are recorded as assumptions in
\path{research/MATH_STATE.json}; their primary passages are archived in
\texttt{literature/}.
\begin{center}\small
\begin{tabular}{p{.18\textwidth}p{.72\textwidth}}
\texttt{MARKMAN-WEIL} &
\cite[Theorem 1.5.1]{Markman}: the whole rational Weil plane of
\emph{every} polarized Weil sixfold of discriminant $-1$ is algebraic.\\
\texttt{MEYER} &
\cite[Corollary 3.3.4]{PS}: a nondegenerate indefinite rational
quadratic form of rank at least five is isotropic.\\
\texttt{HERMITIAN-}\newline\texttt{CONVENTION} &
\cite[Lemma 5.2(2)--(5), author version]{vG}: the homological
polarization formula and determinant convention specified below.\\
\texttt{FFT-SP} &
\cite[Theorem 4.1 and Corollary 4.2]{LZ15}: all-degree
symplectic and ordinary orthogonal invariant spanning by pairs.\\
\texttt{FFT-SO} &
\cite[Theorem 2.1, Remark 2.2, Example 4.1, Corollary 4.2]{LZ17}:
the determinant-isotypic orthogonal component is generated by volume
and pairs, giving the special-orthogonal statement.\\
\texttt{FFT-SL} &
\cite[Section 5, author PDF p. 10]{Sch}: all mixed-degree
special-linear invariant spanning by coevaluation and the two volumes.
\end{tabular}
\end{center}
No CM Hodge-conjecture theorem is an input to (2). The cited
invariant-theory assertions are spanning theorems without a stable-range
restriction; independence of generators is not required.
Markman uses ``complex multiplication by $K$'' for the quadratic action
in his definition of Weil type, not for a toric full Hodge group.
His Theorem 1.5.1 imposes no simplicity or general-period hypothesis.
We import that theorem, not a new verification of his geometric proof.
The Grenoble endpoint of \cite{PS} timed out; the author's Eindhoven
endpoint supplied the stated June 2024 version, dated September 16, 2024,
including the exact corollary. No mathematical conclusion was drawn from
the retrieval failure.

\section{Algebraic operations and rational lifts}

\begin{lemma}\label{lem:descent}
Let $T$ be a rational vector space, $P\subset T_\C$ a complex-linear
subspace, and $t\in T_\kk$. If $\gamma t\in P$ for every
$\gamma\in\Gal(\kk/\Q)$, then
$t\in(T\cap P)\otimes_\Q\kk$.
\end{lemma}
\begin{proof}
Choose a finite Galois field $F$ containing the coefficients of $t$.
For a basis $b_i$ of $F/\Q$ let $b_i^*$ be its trace-dual basis.
The rational coefficient of $b_i$ in $t$ is
$\Tr_{F/\Q}(b_i^*t)$, a sum of scalar multiples of conjugates of $t$.
It lies in $P$. This recovers every coefficient over $\Q$.
The argument does not assert that a scalar automorphism commutes with
complex conjugation.
\end{proof}

Recall that an admissible summand is the image of a $\kk$-projector all
of whose scalar conjugates are Hodge endomorphisms. Admissible maps
preserve the Hodge grading in every scalar conjugate. These are linear
conditions, so Lemma~\ref{lem:descent} expresses the ambient extensions
of such maps as $\kk$-linear combinations of rational Hodge maps.
The degree-one realization used in the reviewed actual-group argument
then realizes them by rational homomorphisms of abelian varieties in the
contravariant direction. Explicitly, dualize to homology, clear lattice
denominators, and use the resulting complex-linear homomorphism of
abelian varieties. Graphs, followed by degree-one projectors, are
algebraic correspondences on the required slots. Isogenies have
cohomological inverses given by degree-normalized pushforward.
All these assertions are cohomological; no Chow-idempotence assertion
is needed.

Pair tensors and contractions also have explicit algebraic realizations.
A rational polarization and the Poincare divisor give a nondegenerate
coevaluation tensor on the ambient degree-one spaces. Any other
all-conjugate $(1,1)$ pair is obtained from it by an admissible map on
one slot. This map is algebraic by the preceding paragraph, so the pair
is algebraic over $\kk$. To contract its inverse pairing, express that
pairing as the polarization pairing composed with an admissible map.
The polarization pairing is cup product with a suitable ample
$(g-1)$st power, followed by integration; hence contraction is algebraic.
This proves the needed abelian-variety instance without an additional
Hodge-conjecture input.

For completeness, a polarization pairs an admissible $W$ perfectly with
its coefficient-conjugate $\bar W$. After complexification, the positive
Hermitian form obtained with the Weil operator detects every nonzero
vector of the Hodge-stable space $W_\C$. Its conjugate arguments lie in
$\bar W_\C$. Conjugating the nonzero determinant preserves perfectness
for every $\gamma$. The rational polarization still pairs complementary
Hodge types on $\gamma W,\gamma\bar W$. Thus
\[
 \epsilon_\gamma(\bar W)=-\epsilon_\gamma(W),\qquad
 \epsilon_\gamma(W)=
 \dim(\gamma W)^{1,0}-\dim(\gamma W)^{0,1}.
 \tag{3}\label{eq:opposite}
\]
This reasoning does not replace $\gamma\bar W$ by $\overline{\gamma W}$.

Write $\delta_r(v_1,\ldots,v_r)=
\sum_{\sigma\in S_r}\operatorname{sgn}(\sigma)
v_{\sigma(1)}\otimes\cdots\otimes v_{\sigma(r)}$ and
$D(W)=\bigwedge^{\dim W}W$ on separate degree-one slots.
For $W=\bigoplus_i W_i$, projecting $\delta_{\dim W}$ onto successive
blocks of lengths $\dim W_i$ gives exactly
$\bigotimes_i\delta_{\dim W_i}$ in concatenated bases.
Conversely, full alternation of this tensor gives
$\left(\prod_i(\dim W_i)!\right)\delta_{\dim W}$.
Both coefficients are nonzero. These are algebraic operations because
the projectors and ordinary permutations are algebraic on these slots;
ordinary and geometric permutations differ by the rational Koszul sign.

If $b:W\otimes W'\to\kk(-1)$ is a perfect pairing of rank $r$,
slotwise contraction gives the exact identity
\[
 b^{\otimes r}\bigl(\delta_r(v)\otimes\delta_r(w)\bigr)
       =r!\det\bigl(b(v_i,w_j)\bigr)\ne0.
 \tag{4}\label{eq:det-contract}
\]
Here the ordinary permutation putting matching slots together is part of
the operation. Alternating the $r$th power of the inverse pair tensor
also supplies $D(W)\otimes D(W')$ nontrivially. In particular every
opposite determinant pair is algebraic. For an alternating self-pair,
the top wedge has a nonzero Pfaffian, so the symplectic volume is already
algebraic. These statements require perfectness, not a particular
choice of determinant vectors or a square root of a scalar.

\begin{lemma}\label{lem:operations}
Supply in (1) is stable under external products, permutations, cup
products, algebraic correspondences, their $\kk$-linear combinations,
and scalar conjugation. Rational membership descends from $\kk$.
Every rational Hodge class in ${\cal C}^{2p}(X)$ has a rational Hodge
lift through finitely many of the defining CM-source maps.
\end{lemma}
\begin{proof}
Products of CM varieties are CM: their joint Hodge group is a subgroup
of a product of tori. External products and composition of cycle
correspondences give the stated operations. Diagonal pullback is cup
product. Expanding a $\kk$-linear combination of maps places its whole
image in the scalar extension of the sum of their rational full images.
The space ${\cal C}^d(X)$ is rational, so its scalar extension is stable
under conjugation and rational membership descends by taking the
quotient by that subspace.

For a target class choose finitely many image maps containing it and
write
\[
 f:\bigoplus_\nu H^{2p-2j_\nu}(B_\nu,\Q)(-j_\nu)
           \longrightarrow H^{2p}(X,\Q).
 \tag{5}\label{eq:lift}
\]
These are polarizable pure rational Hodge structures of the same weight.
Let $K=\ker f$. For a rational polarization on the source, its
restriction to $K$ is nondegenerate: $K$ is stable under the Weil
operator, and the associated positive Hermitian form detects each
nonzero element of $K_\C$. Therefore $K^\perp$ is a rational Hodge
complement, and $f|_{K^\perp}$ is a Hodge isomorphism onto its image.
Its inverse gives a rational Hodge lift. No algebraicity of this
splitting is asserted or needed.
Before twisting, the $\nu$th component has type
$(p-j_\nu,p-j_\nu)$ on $B_\nu$.
\end{proof}

\section{All-degree invariant spanning and one fixed CM list}

Use the reviewed decomposition of $H=H^1(A,\Q)$ into classical blocks.
Fix a finite list ${\cal E}_0$ of embedding lines for the rational
CM isotypic summand $H_0$. This choice precedes every determinant list.
Repetitions retain their chosen lines under scalar conjugation.
Rational toric summands are realized by actual CM isogeny factors
through their degree-one projectors.

The symplectic FFT spans each invariant tensor power by inverse-form
pairs. For the special orthogonal case one must combine the ordinary
orthogonal FFT with the determinant-isotypic assertion of
\cite[Theorem 2.1]{LZ17}. Its harmonic tensor can be chosen to be the
top exterior volume: Example 4.1 verifies harmonicity and nonzero norm.
The result spans by pairs and volume in every degree. For $\SL$ the
mixed-degree FFT spans by identity coevaluations, the volume, and the
dual volume. On Hodge spaces the opposite column is
$\bar W\simeq W^*(-1)$, so the coevaluation is the type-$(1,1)$ pair,
and the opposite volume has positive tensor degree, not the weight of
an untwisted dual. No Tate shift is discarded.

The complex statements descend to $\kk$: in a fixed tensor degree,
invariants of each connected classical group are the common kernel of
the finitely many matrices of its Lie algebra action. Kernels commute
with extension from $\kk$ to $\C$, and the displayed generators are
defined over $\kk$. For a product of groups, decompose the finitely many
slots into their summands and apply the assertions to each factor.
Thus, for any fixed tensor degree, the invariant subspace is spanned by
a \emph{finite} list of generator lines made from pairs, CM embedding
lines, orthogonal volumes, and linear volumes of either orientation.
Each line is pure of some Hodge bidegree after every scalar conjugation.

\begin{lemma}[Simultaneous balance]\label{lem:simultaneous}
Suppose a subspace $U\subset T_\kk$ of a rational pure Hodge structure
of weight $2p$ is spanned by finitely many lines $L_i$, each pure after
every scalar conjugation. Every rational $(p,p)$ vector in $U$ is spanned
by the $L_i$ whose every conjugate has type $(p,p)$.
\end{lemma}
\begin{proof}
For a fixed $\gamma$, group the original lines by the Hodge bidegree
of $\gamma L_i$. Their grouped spans form a direct sum: apply $\gamma$
and use the direct Hodge decomposition. Let $q_\gamma$ retain the
$(p,p)$ group. A rational $(p,p)$ vector $\alpha$ is fixed by
$q_\gamma$, since $\gamma\alpha=\alpha$.
Each $L_i$ is an eigenline of every $q_\gamma$, with eigenvalue zero
or one. It follows on a spanning set that these projections commute,
even when the lines are dependent. Only finitely many distinct
projections occur, since their retained subsets determine them on the
spanning lines. Their product fixes $\alpha$ and has image spanned
by exactly the lines retained by all of them.
\end{proof}

A rational balanced tensor is fixed by its actual rational Hodge group:
its stabilizer is rational and contains the Hodge circle. It is therefore
fixed by the independently acting classical product. Apply the preceding
FFT and Lemma~\ref{lem:simultaneous}. Removing pair factors from the
bidegree count gives precisely
\[
 \sum_j\epsilon_\gamma(W_j)+
 \sum_{\ell\ {\rm in\ the\ chosen\ CM\ list}}\epsilon_\gamma(\ell)=0
       \quad\hbox{for every }\gamma.
 \tag{6}\label{eq:balance}
\]
The list here is the \emph{actual list in that generator}; its entries
are not chosen anew for different $\gamma$.
If every determinant product selected by (6) is supplied, all generator
lines and hence the rational tensor are supplied by
Lemma~\ref{lem:operations}. This proves the unconditional
\emph{supply} version of the manuscript's tensor criterion.
The argument is repeated for each finite degree, with no bound on that
degree or on the number of copies of a block.

\section{Polarization-compatible quadratic determinant supply}

Write $E=\Q(u)$, $u^2=-d$, $d>0$. A rational Riemann form $e$ is
compatible when $e(ux,y)=-e(x,uy)$. Set, on homology,
\[
 h(x,y)=e(x,uy)+u e(x,y).
 \tag{7}\label{eq:hermitian}
\]
It is conjugate-linear first, linear second, Hermitian, and
\[
 e(x,y)=\Tr_{E/\Q}\bigl((2u)^{-1}h(x,y)\bigr).
 \tag{8}\label{eq:trace}
\]
For example $h(x,uy)=-de(x,y)+u e(x,uy)=u h(x,y)$;
the first-variable and conjugate-symmetry identities follow in the
same way from compatibility and alternation. Nondegeneracy follows
from (8). This is exactly \cite[Lemma 5.2]{vG}, the homological
convention used with the discriminant in \cite[Section 1.1]{Markman}.
Compatibility also gives $a^*e=\Nm(a)e$, the polarization condition
in Markman's definition.

For $\sigma(u)=i\sqrt d$, if the $\sigma$-eigenspace of cohomology has
Hodge dimensions $(p,q)$, the signature of (7) is $(p,q)$.
Indeed on the two complex-structure eigenspaces of homology,
$u=\sqrt d I$ and $u=-\sqrt d I$; the signs of $h(x,x)=e(x,ux)$
are therefore the signs claimed by $e(x,Ix)>0$. The summands are
orthogonal because $h$ is $I$-invariant.

There is always a compatible polarization. Starting with $e_0$, take
\[
 e(x,y)=e_0(x,y)+d^{-1}e_0(ux,uy).
 \tag{9}\label{eq:average}
\]
Both summands are positive Riemann forms and direct substitution gives
compatibility. Multiplication by a positive integer clears lattice
denominators. In rank six such multiplication by $N$ changes the
Hermitian determinant by $N^6=\Nm(N^3)$; it does not change its
discriminant class. The same is true of a basis change, which multiplies
the determinant by the norm of the change-of-basis determinant.

\begin{lemma}\label{lem:hyperbolic}
A nondegenerate $E$-Hermitian form of rank $2b$ and signature $(b,b)$
is hyperbolic if and only if $[\det h]=[(-1)^b]$ in
$\Q^\times/\Nm(E^\times)$.
\end{lemma}
\begin{proof}
Given an isotropic $v\ne0$, choose $w$ with $h(v,w)=1$ and replace
$w$ by $w-h(w,w)v/2$. This gives the plane
$\left(\begin{smallmatrix}0&1\\1&0\end{smallmatrix}\right)$.
If a half-dimensional isotropic space exists, projecting it into
the orthogonal complement and iterating gives $b$ such planes.
Their determinant class is $[(-1)^b]$.

Conversely diagonalize $h$ over $E$, with nonzero rational coefficients
$a_i$. The underlying rational quadratic form is
$\sum_i a_i(x_i^2+d y_i^2)$, of dimension $4b$ and signature
$(2b,2b)$. For $b\geq2$, Meyer applies and gives a nonzero isotropic
vector. Split the preceding plane; the complement has the same
signature/determinant condition with $b$ decreased by one.
At rank two write $h=\langle a,b'\rangle$. The determinant condition
gives $\Nm(t)=-ab'$ for some $t\ne0$ in $E$. Then
$(t/a,1)$ is isotropic, since $a\Nm(t/a)+b'=0$.
This finishes the induction. In the applications Meyer is used only
in dimensions twelve and eight; the dimension-four step uses the
explicit norm identity.
\end{proof}

Consequently Markman's discriminant $-1$ at rank six is precisely
the hyperbolic condition for the \emph{actual} compatible form here.
We do not replace that form by an arbitrary form of the same real
signature.

\begin{proposition}[Padding and extraction]\label{prop:padding}
Let $B$ have an $E$-action with compatible signature $(p,q)$,
$n=p+q$. If $1\leq n<6$ and $p,q\leq3$, both $E$-embedding
determinants of $H^1(B,\kk)$ are CM-supplied on separate slots.
No discriminant restriction is needed in this smaller dimension.
If $n=6$ and the actual form is hyperbolic of signature $(3,3)$,
both determinants are algebraic over $\kk$ on separate slots.
\end{proposition}
\begin{proof}
In rank six Markman's theorem supplies the entire rational Weil plane
in $H^6(B,\Q)$. Its two embedding lines are therefore algebraic
over $\kk$. The spectral projectors of $u$ are algebraic on degree one.
To put a determinant on separate slots use addition
$\mu:B^6\to B$ and the algebraic projector
\[
 \pi_1=\prod_{\substack{0\leq r\leq2\dim B\\r\ne1}}
                \frac{[2]^*-2^r}{2-2^r}.
 \tag{10}\label{eq:projector}
\]
Multiplication by $2$ acts by $2^r$ on $H^r(B)$, so
$(\pi_1)^{\boxtimes6}\mu^*(v_1\wedge\cdots\wedge v_6)=\delta_6(v)$.
This is nonzero and algebraic.

Now add $3-p$ $E$-CM elliptic curves with signature $(1,0)$ and
$3-q$ with signature $(0,1)$. Their Hermitian coefficients in
$E$-bases have the prescribed signs and any chosen positive rational
magnitudes: scale their rational polarizations positively.
The actual product $X=B\times\prod_{\nu=1}^{6-n}C_\nu$ has
signature $(3,3)$ and a rational determinant $D_0<0$.
There is at least one added line. Multiply its coefficient by
$-1/D_0>0$. This preserves Riemann positivity, leaves the restricted
form of $B$ unchanged, and makes the total determinant exactly $-1$.
Lemma~\ref{lem:hyperbolic} makes the form hyperbolic. Clearing integral
denominators preserves its discriminant, so Markman applies to this
actual polarized product, at the original period of $B$.

Let $L_\nu=H^1(C_\nu,\kk)_\sigma$. Algebraic block projection of
the supplied total determinant, as calculated before (4), supplies
\[
 D(H^1(B,\kk)_\sigma)\otimes\bigotimes_\nu L_\nu.
 \tag{11}\label{eq:padded-det}
\]
The opposite lines $\bar L_\nu$ are supplied from the \emph{full}
cohomology of the CM product $C=\prod_\nu C_\nu$ by its identity
correspondence. Tensor with these lines and contract the perfect
elliptic pairings. Each contraction is algebraic and nonzero by (4).
Thus (11) gives the desired line, not merely a possibly zero image.
Scalar conjugation supplies the opposite determinant.

One can check the shifts explicitly. The source subspace is
$\bigotimes_\nu\bar L_\nu\subset H^{6-n}(C)$ of Hodge type
$(3-q,3-p)$. The resulting correspondence has shift $j=n-3$:
\[
 H^{6-n}(C)(3-n)\longrightarrow H^n(B^n).
 \tag{12}\label{eq:padding-shift}
\]
After twisting its indicated type is exactly $(p,q)$.
The correspondence cycles on $C\times B^n$ have codimension
$(6-n)+(n-3)=3$ in every smaller rank.
Negative shifts for small $n$ are legitimate in (1).
For $n=6$ the source is a point; $D(0)$ is also supplied by a point.
\end{proof}

\section{Exhaustive special-period balance}

Work up to isogeny as $A=\prod B^{r_B}$, transporting the given action.
Pullback preserves the rational Hermitian isometry class, and the
degree-normalized inverse correspondence preserves supply.
The isotypic summands are orthogonal for any compatible polarization:
the polarization has no cross homomorphisms between distinct simple
isogeny types and their duals. Their projectors commute with $E$.
In the sixfold case their signatures therefore add to $(3,3)$.
A proper summand need \emph{not} itself be hyperbolic.

Fix the CM line decomposition as specified below before testing any
list in (6). For each linear representative $W_i$, let
$P_i,Q_i\geq0$ be the counts of $D(W_i),D(\bar W_i)$ and set
$n_i=P_i-Q_i\in\Z$. Orthogonal columns have zero imbalance.
Write $c_\gamma$ for the sum from the fixed CM list. Balance implies
\[
 \sum_i n_i\epsilon_\gamma(W_i)+c_\gamma=0
                    \quad\hbox{for every }\gamma.
 \tag{13}\label{eq:net}
\]
The $\min(P_i,Q_i)$ opposite pairs are already algebraic.
If several $n_i$ are equal to one integer $t$, their remaining line is
the $t$th power of the product of those determinants for $t\geq0$,
or the $(-t)$th power of its opposite for $t<0$.
There is no division by an integer exponent or extraction of a root.

\subsection{Blocks with direct supply or forced cancellation}

Types I and II require only pair tensors. A Type III factor has rank
four or six; its $E$-eigenspace is an admissible copy of its orthogonal
column. Rank four has signature $(2,2)$ and is supplied by
Proposition~\ref{prop:padding} for its \emph{arbitrary} restricted form.
Rank six exhausts $A$ and uses the stipulated hyperbolic form.
Admissible copy maps transfer these supplies to the representative.

For Type IV with quadratic center equal to $E$, a rank-three or
rank-four column is mixed and both signature entries are at most three,
so its determinant is individually supplied.
In dimension at most five a rank-five factor exhausts $A$ and has
imbalance $2p-5\ne0$. There is no CM remainder or other block.
Equation (13) therefore gives $n=0$, including signatures $(1,4)$
and $(4,1)$ which the padding proposition does \emph{not} cover.

In the split sixfold, a multiplicity-one factor of rank at most five
has both signature entries at most three by orthogonal addition, so
padding supplies it. A rank-six factor exhausts $A$ and uses the
given hyperbolic form. The only possible repeated non-CM factor in
this case is a threefold $B^2$. For the central $K=E$ action on $B$
the signature is $(1,2)$ or $(2,1)$, hence padding applies on $B$.
The given, possibly noncentral, action on $B^2$ merely reassigns its
two opposite column types among matrix copies. Their admissible
identifications preserve the supply; a product polarization on $B^2$
is not substituted for the given sixfold polarization.

If $E$ does not embed in the Type IV center $K$, the reviewed list
forces a quadratic-center threefold repeated twice, exhausting $A$.
Its single $K$-column has rank three and nonzero imbalance.
There is no CM remainder, so (13) again gives zero net exponent.
This case does not occur below dimension six.

These arguments are unaffected by other symplectic blocks, by arbitrary
CM remainders in the individually supplied cases, or by two distinct
quadratic-center threefolds in dimension six.

\subsection{Quartic-center fourfold and its fixed CM remainder}

The remaining proper higher-center factor has $(f,a,g)=(2,2,4)$.
Take representatives $W_1,W_2$ above one embedding of $E$.
Exactly one is mixed and the other pure, so the imbalance vector is,
up to order, $(0,2)$ or $(0,-2)$.

Below dimension six the remainder has dimension zero or one.
In dimension one its inherited unital $E$-action realizes an $E$-CM
elliptic curve. In a split sixfold the restricted signature of $B$ is
$(1,3)$ or $(3,1)$; the remainder has respectively signature $(2,0)$
or $(0,2)$. On its two-dimensional $E$-homology the complex structure
is the same $E$-scalar everywhere. Choose an $E$-basis: each basis
line is a rational Hodge substructure of elliptic CM type.
Degree-one realization makes the remainder isogenous to a square of
an $E$-CM elliptic curve. This is a geometric consequence of the
signature, not an assumption that all CM factors are elliptic.

In both families choose the embedding lines of these $E$-CM factors
as the fixed CM decomposition. Every $\gamma\in\Gal(\kk/E)$ fixes
each of these lines and hence their contribution $c_\gamma=c$,
regardless of the number of repeated entries in the list.
An element of this Galois group exchanges the two embeddings of $K/E$.
Writing $d_i=\epsilon_1(W_i)$, (13) includes
\[
 d_1n_1+d_2n_2+c=0,\qquad
 d_2n_1+d_1n_2+c=0.
 \tag{14}\label{eq:quartic}
\]
As $d_1-d_2=\pm2$, subtraction gives $n_1=n_2=t\in\Z$.
The integer kernel in variables $(n_1,n_2,c)$ is exactly
$\Z(1,1,-d_1-d_2)$; its first coordinate is one, so no saturation
or divisibility issue is hidden.

Algebraic direct-sum determinant maps identify the basic product
$D(W_1)D(W_2)$ with the full $E$-determinant of $B$.
Its signature is $(1,3)$ or $(3,1)$, so padding supplies it.
Integer cancellation then supplies every determinant list in (13).
If the remainder is zero, extra balance may force $t=0$; supplying
the larger set just obtained is harmless.

\subsection{Higher-center sixfolds}

Here $B=A$ and there is no CM remainder. The given action is central,
since the multiplicity is one. Representatives are taken above one
$E$-embedding. Scalar automorphisms fixing $E$ permute these embeddings;
their action is transitive. Primitive matrix-copy changes, when present
in the general construction, are admissible and do not change the
Hodge multiplicities.

If $(f,a)=(2,3)$, the two holomorphic multiplicities sum to three.
The possibilities $(0,3),(3,0)$ have every embedding pure and are
CM, contrary to this case. Thus the vector is $(1,2)$ or $(2,1)$,
with imbalances $(-1,1)$ or $(1,-1)$. Balance gives $n_1=n_2$.
The integer kernel is $\Z(1,1)$.

If $(f,a)=(3,2)$, the multiplicities lie in $\{0,1,2\}$ and sum
to three. The all-mixed vector $(1,1,1)$ is excluded by the reviewed
rank-two actual-commutant argument. The remaining possibilities are
the permutations of $(0,1,2)$. The transitive action on three
embeddings contains a three-cycle. In cyclic order the necessary
balance equations have rows which are cyclic permutations of
$(-2,0,2)$, up to a common relabeling and reversal. For example
\[
 M=\begin{pmatrix}-2&0&2\\2&-2&0\\0&2&-2\end{pmatrix},
       \qquad
 \ker_\Q M=\Q(1,1,1),\quad
 \ker_\Z M=\Z(1,1,1).
 \tag{15}\label{eq:cyclic}
\]
Thus $n_1=n_2=n_3=t\in\Z$.

In either case $\bigotimes_{i=1}^fD(W_i)$ is identified nontrivially
and algebraically with the full $E$-determinant of $A$.
The actual hyperbolic form supplies this determinant; it is not
necessary to supply each column determinant separately.
Opposite cancellation and the integer $t$ finish the argument.
Further scalar-conjugate equations can only restrict these necessary
kernels, never invalidate the supply constructed for them.

\begin{proposition}\label{prop:balance}
For each family in Theorem~\ref{thm:main}, the geometry permits one
fixed CM-line decomposition for which every determinant list satisfying
(6), with arbitrary repetitions, is CM-supplied. Thus condition (d)
of the manuscript's tensor criterion holds.
\end{proposition}
\begin{proof}
The preceding cases exhaust the reviewed covering list. A determinant
block has cost at least three in abelian dimension, an orthogonal
block at least four, and each higher-center exceptional case at least
four. A quartic-center fourfold is consequently the only nonsaturated
proper case that can interact with a CM remainder, which was treated
with a fixed decomposition above. Non-embedded quadratic actions and
higher-center sixfolds exhaust the dimension budget; the odd rank-five
case below six likewise exhausts it. Other blocks are individually
supplied. All exponent arguments were over $\Z$, without bounds.
\end{proof}

\section{Completion on powers and the conditional implication}

\begin{proof}[Proof of Theorem~\ref{thm:main}]
Apply the reviewed group theorem at the actual period.
The FFT reduces invariant tensors to generator lines.
Lemma~\ref{lem:simultaneous} retains only those balanced in all conjugates.
Proposition~\ref{prop:balance} supplies their determinant factors.
The CM embedding lines come from actual CM factors, and pair tensors
are algebraic. Lemma~\ref{lem:operations} descends this supply rationally.

On $A^m$ first cohomology is $H^{\oplus m}$. Label all slots by their
copies. This just repeats the original blocks and CM lines; the
unbounded repetitions already allowed in (6) and
Proposition~\ref{prop:balance} apply. For a class in
$H^{2p}(A^m,\Q)=\bigwedge^{2p}H^1(A^m,\Q)$, normalized
antisymmetrization is a rational Hodge lift to the ordinary tensor power.
The proved tensor result supplies that lift on separate slots.
Diagonal pullback, equivalently cup product of those slots, returns
the original exterior class. Supply is preserved by this algebraic map.
The case $p=0$ is supplied by a point, and nonexistent degrees are zero.
This proves (2) for every $m,p$.

The rational Hodge lift is furnished by (5).
Only for the final conditional implication do we assume the CM Hodge
theorem: it algebraizes each untwisted rational Hodge component of that
lift on its CM source. Applying the given rational algebraic
correspondences and adding then algebraizes the target.
\end{proof}

This conclusion is not a claim that the full cohomology of $A$ is
CM-supplied, or that its derived group is trivial. Nor does it assert
saturation under arbitrary additional CM companions. For example the
rank-five cancellation argument used the absence of a CM remainder.
The general saturated extension in the manuscript is a separate
conditional statement and is not needed here.

\section{Exact certificate and what it proves}

The read-only command
\[
 \texttt{python3 -B research/determinant\_supply\_audit/verify.py}
\]
recomputes \texttt{exact\_certificate.json}. Its only mathematical
dependency beyond the Python standard library is SymPy, used in exact
rational and symbolic arithmetic.
It does not call a network service or trust precomputed matrix ranks.
The certificate checks:
\begin{enumerate}
\item The generic averaging, Hermitian, trace and hyperbolic-plane
identities, the determinant-normalizing positive scaling formula,
and the degrees and Hodge shifts in (12).
\item Nonzero separate-slot extraction for every composition of ranks
one through six, the generic polynomial contraction identity (4),
and the Pfaffian sign identities in symplectic ranks two, four and six.
\item The degree-one projector at every cohomological degree for
abelian dimensions one through six.
\item Primitive integer kernels for every orientation/permutation of
the exceptional balance vectors, including the fixed CM contribution
in (14); universal symbolic nonnegative monoid reconstruction for
both signs of the common integer exponent.
\item All determinant-bearing configurations within the dimension
budget and permitted signatures of the reviewed covering list.
Unspecified residual dimensions overapproximate the possible
pair-only or CM factors. Every exceptional residual signature used
in the proof is checked explicitly.
\end{enumerate}
This finite coverage is exhaustive because dimensions and column ranks
are bounded, \emph{not} because tensor degree or determinant exponents
are bounded. The proof above supplies the all-degree invariant theorem,
fixed-list quantifier, geometric realization of the residual CM factors,
and arbitrary-integer cancellation. The program does not decide which
numerical configurations are realized by an abelian variety, prove
Markman's theorem or Meyer anew, or verify algebraicity of arbitrary
CM Hodge classes.

The saved Jacobian request concerns only the $3$-by-$3$ rational matrix
(15). Its current typed contract was inspected after discovery, and the
approved Argus adapter was then run once. The operation id, payload,
server/protocol metadata, typed output and adapter exit status are
preserved beside the certificate. The standalone checker recomputes
the nullspace independently and verifies the returned basis. A timeout,
transport error or incomplete result would fail this replay rather than
count as mathematical evidence.

Negative controls show why weaker tests do not suffice. The identity
equation alone for $(-2,0,2)$ admits $(0,1,0)$, whereas the cyclic
system excludes it. Allowing an independent CM contribution for each
conjugate of a quartic-center factor also permits unequal exponents.
These are counterexamples to weakened \emph{linear systems}, not
hypothesis-valid abelian-variety counterexamples. The geometry above
is exactly what rules them out. We do not promote them to a refutation.

This extends the reviewed group result by proving the requested supply
reduction. It does not establish publication priority. Independent review
of the proof and citations is still required, as is any assessment of
novelty and significance. The Kuga--Satake, geometric Weil-plane proof
and Milne branches, and the companion CM theorem, are untouched.

\begin{thebibliography}{9}
\bibitem{Markman}
E. Markman, \emph{Cycles on abelian $2n$-folds of Weil type from secant
sheaves on abelian $n$-folds}, arXiv:2502.03415v2, August 24, 2026.
\url{https://arxiv.org/html/2502.03415v2}.
\bibitem{PS}
C. Peters and H. Sterk, \emph{Symmetric and Quadratic Forms, with
Applications to Coding Theory, Algebraic Geometry and Topology},
author monograph, Version June 2024, dated September 16, 2024,
Corollary 3.3.4, p. 77.
\url{https://cpeters1.win.tue.nl/Books/QuadraticForms/QuadForms.pdf}.
\bibitem{vG}
B. van Geemen, \emph{An introduction to the Hodge conjecture for abelian
varieties}, LNM 1594 (1994), 233--252,
doi:10.1007/978-3-540-49046-3\_5.
Lemma locators refer to \url{https://sites.unimi.it/vangeemen/cime.pdf}.
\bibitem{LZ15}
G. I. Lehrer and R. B. Zhang, \emph{The Brauer category and invariant
theory}, JEMS 17 (2015), 2311--2351, doi:10.4171/JEMS/558.
\bibitem{LZ17}
G. I. Lehrer and R. B. Zhang, \emph{Invariants of the special orthogonal
group and an enhanced Brauer category}, Enseign. Math. 63 (2017),
181--200, doi:10.4171/LEM/63-1/2-6.
\bibitem{Sch}
A. Schrijver, \emph{Tensor subalgebras and first fundamental theorems
in invariant theory}, J. Algebra 319 (2008), 1305--1319,
doi:10.1016/j.jalgebra.2007.10.039.
Section/page locators refer to the 19-page author PDF:
\url{https://homepages.cwi.nl/~lex/files/tensfft_ja.pdf}.
\end{thebibliography}
\end{document}
